\expandafter\def\csname HistoricalNoteText20\endcsname{The preceding article already gives some examples of the considerable disagreement between different observers at the same place. A few more may be cited here in which the differences are far too great to be attributed with any probability to regular annual change. At Valparaiso the inclination in 1829, according to \textsc{King}, was $-40^0\,11'$, and in 1835, according to \textsc{Fitz Roy}, it was $-38^0\,3'$. On the island of Mauritius the intensity in 1818, according to \textsc{Freycinet}, was $1{,}096$, and in 1836, according to \textsc{Fitz Roy}, it was $1{,}192$. The difference is greater still at Otaheite, where in 1830 \textsc{Erman} found the intensity $=1{,}172$, whereas in 1835 \textsc{Fitz Roy} found it $=1{,}017$. This last difference, at a place of the highest importance for future improvements to the elements, is considerably greater than the largest difference among all 86 of our comparisons of calculated with observed intensities.}
\SourceLeaf{166}
\Article{30}
In judging the differences between calculation and observation shown in the preceding table, we have to take two things into account. Almost all the observations are affected by errors in making the measurements and by incidental anomalies in the magnetic force itself, and the observations do not all apply to the same year\HistoricalNote{20}. On the other hand, our formulas contain only terms up to the fourth order, whereas the later terms may still be quite appreciable. The agreement between calculation and observation thus seems to meet all reasonable expectations for a first attempt. Our expression for $\frac VR$ may therefore be regarded as close to the truth, at least in its larger terms. It therefore seemed worth showing graphically how the numerical values of $\frac VR$ vary. This has been done in a map drawn by Dr. \textsc{Goldschmidt} in three parts. The first, on \textsc{Mercator}'s projection, shows the entire belt around the Earth between $70^0$ north and $70^0$ south latitude. The other two, on stereographic projections, show the polar regions as far as $65^0$ latitude. The corrections and additions to the expression for $\frac VR$ that a new calculation based on better data will require will undoubtedly still cause considerable shifts in this system of lines, especially at high southern latitudes. But an essential change in the overall pattern is inconceivable without changes in the expression for $\frac VR$ so large that agreement with the existing observations would be lost. We are thus led to the important result that the system of
\SourceLeaf{167}
lines of equal values of $V$ on the Earth's surface does indeed fall under the simplest type described above in Art. 11, so that there are \textit{only two magnetic poles} on the Earth if we set aside the case of a local exception mentioned in Article 13, whose occurrence or nonoccurrence must remain undecided for now. An exact calculation from our elements gives the positions of these two poles

1) at $73^0\,35'$ north latitude, $264^0\,21'$ longitude east of Greenwich, with total intensity $=1{,}701$ (in the usual unit).

2) at $72^0\,35'$ south latitude, $152^0\,30'$ longitude, with total intensity $=2{,}253$.

At the first point, $\frac VR$ has its greatest value $=+895{,}86$, and at the second its smallest $=-1030{,}24$.

According to \textsc{Ross}'s observation, the northern magnetic pole is $3^0\,30'$ farther south than in our calculation. As our comparison table shows, the calculation gives a direction of the magnetic force at that place that is in error by $1^0\,12'$. A considerably greater displacement is to be expected at the southern magnetic pole. Hobarttown is the observation site nearest that pole. There our calculated inclination, disregarding the sign, is $3^0\,38'$ too small, insofar as the observation can be relied on. The actual southern magnetic pole therefore probably lies considerably farther north than our calculation indicates and might be sought in the vicinity of $66^0$ latitude and $146^0$ longitude.

\Article{31}
Some significance may be attached to the two points on the Earth's surface where the horizontal force vanishes and which are called the magnetic poles, because of their relation to the overall pattern of the horizontal force on the Earth's surface. But we must take care in extending that significance any further. In particular, the chord joining the two points has no significance, and it would be quite wrong to single out \textit{this} straight line by calling it the Earth's \textit{magnetic axis}. The only way to give the concept of a body's magnetic axis a
\SourceLeaf{168}
generally valid meaning is the one laid down in Article 5 of the \textit{Intensitas vis magneticae}. There it means a straight line with respect to which the moment of the free magnetism contained in the body is a maximum. To determine the position of the Earth's magnetic axis in this sense, and at the same time the moment of terrestrial magnetism with respect to it, we need only know the first-order terms of $V$, as already noted above in Article 17. Our elements in Article 26 give
\[
P'=+925{,}782\cos u+89{,}024\sin u\cos\lambda-178{,}744\sin u\sin\lambda
\]
so that $-925{,}782R^3$, $-89{,}024R^3$, $+178{,}744R^3$ are the moments of terrestrial magnetism with respect to the Earth's axis and the two radii at longitudes $0$ and $90^0$. The direction meant for the Earth's axis is toward the north pole. The negative sign of the corresponding moment shows that the magnetic axis makes an obtuse angle with it, so that its magnetic north pole is turned toward the south. From this the direction of the magnetic axis is found to be parallel to the diameter of the Earth running from $77^0\,50'$ N. latitude, $296^0\,29'$ longitude to $77^0\,50'$ S. latitude, $116^0\,29'$ longitude, and the magnetic moment with respect to it is $=947{,}08R^3$. For the latter, we must remember that our elements use a unit of intensity that is one thousandth of the usual one. To convert to the absolute unit established in the \textit{Intensitas vis magneticae}, we note that in that work the horizontal intensity at Göttingen, on 19 July 1834, was found to be $=1{,}7748$. With the inclination $68^0\,1'$ this gives a total intensity $=4{,}7414$, while in the unit used above it is taken to be $=1357$. The conversion factor is therefore $=0{,}0034941$, and so the Earth's magnetic moment in the absolute unit is
\[
=3{,}3092R^3
\]
Since the millimeter is taken as the unit of length in this absolute unit of terrestrial magnetic force, $R$ too must be expressed in millimeters. As the ellipticity of the Earth is being left out of account here in any case, it suffices to take $R$ as the radius of a circle with circumference $40000$ million millimeters. The magnetic moment above is thus expressed by
\expandafter\def\csname HistoricalNoteText21\endcsname{Insofar, that is, as we are not entitled to assume that the magnetic axes of all the magnetized parts of the Earth are strictly parallel. The more these axes depart from parallel alignment, the stronger the average magnetization of the parts must be to produce the same total magnetic moment.}
\SourceLeaf{169}
a number whose logarithm is $=29{,}93136$, or by $853800$ quadrillion on the long scale. In the same absolute unit, the experiments made in 1832 gave the magnetic moment of a one-pound bar magnet as $=100877000$ (\textit{Intensitas} Art. 21). The Earth's magnetic moment is thus $8464$ long-scale trillion times as great. It would therefore take $8464$ trillion such bar magnets, with parallel magnetic axes, to replace the magnetic effect of the Earth in the space outside it. Distributed uniformly throughout the volume of the Earth, this comes to nearly eight bars (more exactly $7{,}831$) per cubic meter. Stated this way, the result keeps its meaning even if we were to ascribe terrestrial magnetism solely to persistent galvanic currents in the Earth rather than consider the Earth an actual magnet. But if we consider the Earth an actual magnet, we are obliged to ascribe to every part of it one eighth of a cubic meter in volume, \textit{on average} at least\HistoricalNote{21}, a magnetization as strong as that of the bar magnet. This is a result that will probably surprise physicists.

